\documentclass[11pt,a4paper]{article}
\usepackage[margin=2.2cm]{geometry}
\usepackage{graphicx}
\usepackage{booktabs}
\usepackage{amsmath}
\usepackage{xcolor}
\usepackage[hidelinks]{hyperref}
\usepackage{caption}
\usepackage{float}
\captionsetup{font=small,labelfont=bf}
\setlength{\parskip}{0.45em}
\setlength{\parindent}{0pt}
\setlength{\emergencystretch}{2em}

\title{Text Summarization with Transformers:\\A Small Length-Controlled BART for the Browser}
\author{Manas Pandya\\[0.2em]\small NLP \& CV Lab, Task 03}
\date{November 2026}

\begin{document}
\maketitle

\begin{center}
\small
Live demo: \url{https://nlp-cv-lab.vercel.app/p/summarization}\\
Code: \url{https://github.com/manasrpandya/NLP_CV_Lab/tree/main/projects/03-summarization}\\
Models: \url{https://huggingface.co/manasrpandya/cnndm-summ-onnx} (browser),
\url{https://huggingface.co/manasrpandya/cnndm-summ} (PyTorch)
\end{center}

\section{Aim}

The task is to summarize text with transformer models. The usual quick answer is to run
\texttt{facebook/bart-large-cnn} through the Hugging Face pipeline, print one summary and maybe a ROUGE
score. I wanted to know a bit more than that, so the plan was a ladder of methods and a proper evaluation
of each:

\begin{enumerate}
\itemsep0em
\item[(a)] extractive baselines: Lead-3 and TextRank,
\item[(b)] pretrained abstractive models used as they are: BART-large-CNN, PEGASUS-large and
DistilBART-CNN-12-6,
\item[(c)] my own model, a BART-base fine-tuned on CNN/DailyMail that is small enough to run in a browser
tab (a 137~MB download after quantization),
\item[(d)] length control: a prefix (\texttt{summarize short:}, \texttt{medium:} or \texttt{long:}) so that
the user can choose how long the summary is, and a check that the output length really follows it,
\item[(e)] sequence-level knowledge distillation from BART-large-CNN, to see if it helps a small model.
\end{enumerate}

Every method is scored with ROUGE, BERTScore, abstractiveness (share of novel bigrams) and a factual
consistency score from an NLI model, and I read twenty summaries by hand to count hallucinations. The
trained model is exported to ONNX and runs on the demo page, which can also check the summary against the
source with a small NLI model, in the browser as well.

\section{Background}

\textbf{Extractive and abstractive.} Extractive summarizers pick sentences from the source. On news they
are strong, because journalists put the key facts at the start: simply taking the first three sentences
(Lead-3) is a baseline that many neural models only just beat \cite{see17}. TextRank \cite{mihalcea04}
scores sentences by running PageRank on a graph where sentences are nodes and edges are weighted by
similarity. Abstractive summarizers write new text. Modern ones are encoder-decoder transformers
pretrained with a denoising or gap-sentence objective and then fine-tuned: BART \cite{lewis20} reconstructs
corrupted text, PEGASUS \cite{zhang20} masks whole sentences and generates them, and T5 \cite{raffel20}
casts everything as text to text. Shleifer and Rush \cite{shleifer20} compare ways to shrink such models
and the ``shrink and fine-tune'' DistilBART models come from that work.

\textbf{Evaluation.} ROUGE \cite{lin04} counts n-gram overlap with a human summary. It is cheap and standard,
but it rewards copying and says nothing about whether the summary is true. BERTScore \cite{zhang20bs}
compares contextual embeddings and is more forgiving of paraphrase. Abstractive models are known to
hallucinate: Maynez et al.\ \cite{maynez20} found that a large share of the output of models on XSum
contains content that is not in the source. A common automatic check is to ask an NLI model whether the
source entails each summary sentence; SummaC \cite{laban22} does this for every pair of summary and source
sentences and aggregates. I use its zero-shot variant.

\textbf{Distillation and length.} Sequence-level distillation \cite{kim16} trains the small model on the
outputs of a big one instead of on the human references. Length control can be done by giving the length
class as an extra input token, as in Fan et al.\ \cite{fan18}, which is what I do with a text prefix.

\section{Data}

\textbf{CNN/DailyMail 3.0.0} \cite{hermann15,see17}: 287k training articles with bullet-point highlights as
the summary, 13k validation and 11.5k test. Articles average about 780 words and reference summaries about
52. I train on a random 40,000 article subset of the training set for one epoch, which fits the Colab
budget. All models are scored on the same fixed random subset of 2,000 test articles (NumPy seed 0, the
indices are recomputed in \texttt{summ.data.subset}). The faithfulness score is the slowest to compute, so
it uses the first 500 of those 2,000. The human reference summaries are the highlights with sentences joined
by newlines.

\textbf{XSum} \cite{narayan18} as the second set: BBC articles with a single abstractive sentence as the
summary, so it is the opposite of CNN/DailyMail, where copying works well. It loaded without trouble, so I
did not need the SAMSum fallback. I fine-tune BART-base on 20,000 articles and score on a fixed random 1,000
test articles.

\section{Method}

\subsection{Extractive baselines}

Lead-3 returns the first three sentences. TextRank builds a TF-IDF cosine similarity graph between sentences,
runs PageRank (damping 0.85) and returns the top three sentences in source order. Sentences are split with
a small regex splitter (\texttt{summ.text.sents}) that knows a few abbreviations; the same rule is
implemented in JavaScript for the ``quick'' mode of the demo.

\subsection{Pretrained references}

BART-large-CNN, PEGASUS-large (\texttt{google/pegasus-cnn\_dailymail}) and DistilBART-CNN-12-6 are used with
the generation settings in their own hub configs (beam search with 4 beams, length penalties and minimum
lengths as released) and sources cut at 1024 tokens. PEGASUS gives empty output in fp16, a known issue, so it
runs in fp32; the other two run in fp16.

\subsection{My model}

\textbf{Choice of model.} The browser limit I set was about 150~MB after int8 quantization. T5-small would
fit easily, but T5-family models overflow in fp16 and the Colab T4 has no bf16, so they have to be trained in
fp32, which is several times slower. BART trains fine in fp16, so I took \texttt{facebook/bart-base}
(6+6 layers, $d=768$, 139M parameters). Its problem is size: exported naively, the int8 encoder is 79~MB and the
merged decoder 129~MB. The embedding matrix ($50{,}265\times768$) appears three times, in the encoder, in the
decoder and in the output layer, so I trimmed the vocabulary (Section~\ref{sec:trim}).

\textbf{Training.} One epoch over 40,000 articles, batch size 16, AdamW with learning rate $5\cdot10^{-5}$,
5\% warm-up and linear decay, label smoothing 0.1, fp16 autocast with a gradient scaler. Inputs are cut at
512 tokens, targets at 128. I keep the epoch whose ROUGE-1 on 200 validation articles is best, which with a
single epoch is the last one. Training took about 20 minutes on a T4.

\textbf{Length prefix.} Every training input starts with \texttt{summarize short: },
\texttt{summarize medium: } or \texttt{summarize long: }. The bucket is taken from the number of words in the
human reference using the tertiles of the training set (41 and 55 words), so each bucket holds a third of the
data. At test time the user picks the prefix. There is no change to the model or the loss.

\textbf{Sequence-level distillation.} For 8,000 training articles I generated BART-large-CNN summaries (beam 2)
and trained a second BART-base on them, next to a control trained on the human summaries of the same 8,000
articles. The bucket for the prefix still comes from the human summary length.

\subsection{Metrics}

\textbf{ROUGE-1/2/Lsum} F1 from the \texttt{rouge\_score} package with stemming; for Lsum each summary is
split into sentences, one per line. 95\% intervals are bootstrap intervals over articles.
\textbf{BERTScore} F1 with \texttt{roberta-large}. \textbf{Novel bigrams} is the share of summary bigrams
(lowercased word tokens) that do not occur in the article. \textbf{Faithfulness} follows SummaC-ZS: an NLI
model (\texttt{cross-encoder/nli-deberta-v3-xsmall}, the same one the demo runs in the browser) scores
every (source sentence, summary sentence) pair, each summary sentence takes its maximum entailment
probability over source sentences, and the summary score is the mean over its sentences. I report it as a
percentage. The original SummaC-ZS uses entailment minus contradiction; I use the plain entailment
probability. For the human references the same score is computed to show what a ``normal'' value is.
\textbf{Length} is the mean number of words, and \textbf{speed} is CPU time per summary with 4 threads on my
laptop, including tokenization, on the first 20 articles of the subset (5 for the large PyTorch models).

\subsection{Getting it into the browser}
\label{sec:trim}

\textbf{Vocabulary trimming.} I counted the BPE token types that occur in the 40,000 training articles and
their summaries (47,819 of 50,265 types appear at least once, mostly rare names) and kept the 18,000 most
frequent ones, all single-byte tokens and the merge ancestors needed to build them. Embeddings, the output
layer and its bias were sliced to those rows, and the tokenizer was rebuilt with the same ids order, so the
special tokens keep ids 0--3. Text that uses a dropped token is split into smaller pieces by BPE. On
test articles this makes the input about 6\% longer in tokens.

\textbf{Export.} Optimum exports the encoder and a decoder with and without the key/value cache. The decoders are
quantized one by one (dynamic 8-bit for matrix multiplications and the embedding gather) and merged
afterwards into one graph with an \texttt{If} node, which is the layout transformers.js expects. Quantizing the
merged graph directly left the decoder in fp32, because the quantizer does not go into the \texttt{If}
branches. Weights are stored unsigned: signed int8 weights lost about 2.4 ROUGE-1 on a 100 article test
and unsigned ones about 1.1, and per-channel quantization gave garbage in this graph.

\textbf{Demo.} The page runs the model in a web worker through transformers.js on the WASM backend with greedy
decoding (no repeated trigrams, at most 140 new tokens) and streams tokens as they are produced. A second small
model runs on request for the fact check: the ONNX build of the same NLI model, 8-bit. The ``quick'' mode is
TextRank in plain JavaScript.

